\subsection{NN-G031: a competitive closed table must have gluing mass at least eight}
\paragraph{Statement.}
Every effective associative symmetric closed table of the specified trace
with at most $3352$ cancelling pairs has $W(4)\ge8$.
This permits arbitrary cell degrees and contents and arbitrary
rational ratios and auxiliary parameters, subject only to
integrality and effectivity of the actual coefficients.
In particular the entire $W=D+2$ branch in this presentation is excluded, without
enumerating support locations or assuming a degree bound.

\paragraph{Mass notation and immediate bounds.}
Put $\psi=D(D-3)=q^2+q+q^{-1}+q^{-2}$, the absolute
cell trace, and use lower-case letters for evaluation at $D=4$.
For $A_i=\delta_i\Lambda_i$, write $a_i=A_i(4)$.
All these masses are positive even integers. The trace imposes
\[
 d_i\ge4,\quad u_i\ge4,\quad v_i\ge16,\qquad
 a_1\ge8,\quad a_2\ge16.
\]
The nonunit mass, with the unit contributions cancelled from
the pair count, is
\[
 T=s+20(g_1+g_2)+200w+20(u_1+u_2+v_1+v_2)
       +26(d_1+d_2)+14(a_1+a_2).                    \tag{G031.1}
\]
The proposed score cap is exactly $T\le11128$.
Associativity gives
\[
 ws=g_1g_2,\qquad u_i g_j=w(g_i-4d_i)\quad(i\ne j).
                                                               \tag{G031.2}
\]
Consequently $u_1u_2<w^2$. Since $u_i\ge4$, $w=4$
is impossible. Parity and domination give $w\ge4$ even,
so it remains to exclude $w=6$.

An effective symmetric polynomial of mass six dominating
$D=q+2+q^{-1}$ and congruent to it modulo two is uniquely
$W=D+2$: its surplus mass two must be central.
Similarly $u_i=4$ forces $U_i=D$, and $u_i=6$ forces
$U_i=D+2=W$.

\paragraph{A useful divisibility consequence of $U_i=D$.}
For this $W$, the identities are
\[
 U_1G_2=W(G_1-D\delta_1),\qquad
 U_2G_1=W(G_2-D\delta_2).                           \tag{G031.3}
\]
If $U_1=D$, reduction of the first at $D=0$ implies
$D\mid G_1$, and coprimality of the monic polynomials
$D,W$ implies $W\mid G_2$. Write $G_1=DR$, $G_2=WT$;
all quotients are integral. The second identity at $D=0$
then implies $T(0)=0$, because $W(0)^2=4$.
Thus
\[
 G_2=DW C_2,\qquad G_1=D\delta_1+D^2C_2.
\]
The parity $G_1\equiv0$ and
$\delta_1\equiv D(D-1)$ in $\mathbb F_2[D]$ gives
$C_2\equiv D-1$ after cancellation of $D^2$.
Therefore $g_2/24$ is an odd integer.
The same argument with indices interchanged shows that
$U_2=D$ implies
\[
 G_1=DW C_1,\qquad G_2=D\delta_2+D^2C_1,
 \qquad C_1\equiv0\pmod2,                          \tag{G031.4}
\]
where the last parity follows from $G_1\equiv0$.
Hence $48\mid g_1$. None of these quotient polynomials
is asserted to be effective.

\paragraph{The five possible $u$ pairs.}
Since $u_1,u_2$ are even and at least four, their product
less than $36$ leaves only
\[
 (4,4),\ (4,6),\ (6,4),\ (4,8),\ (8,4).
\]
Solving (G031.2) at $w=6$ gives
\[
 g_1=\frac{24(6d_1+u_1d_2)}{36-u_1u_2},\qquad
 g_2=\frac{24(u_2d_1+6d_2)}{36-u_1u_2}.             \tag{G031.5}
\]
In either pair containing eight, $g_1+g_2\ge576$,
so $20(g_1+g_2)>11128$ already.

For $(u_1,u_2)=(6,4)$, the actual identities with
$U_1=W,U_2=D$ yield
\[
 G_1=DW K,\qquad K=(\delta_1+\delta_2)/2.
\]
The polynomial $K$ is integral and effective, with
$K\ge\psi$ coefficientwise. Since $G_1\equiv0$ and
$DW\equiv D^2$ modulo two, $K$ has even coefficients.
Its coefficients at $q^{\pm1},q^{\pm2}$ are therefore
at least two. Hence $t=d_1+d_2=2K(4)\ge16$.
Formula (G031.5) gives $g_1=12t$ and
$g_2=8t+4d_2\ge8t+16$. Consequently
\[
 T\ge s+20(g_1+g_2)+1200
    \ge16t^2+432t+1520
    \ge12528>11128.
\]

For $(u_1,u_2)=(4,6)$, the identities instead give
$G_2=DW K$ with the same effective
$K=(\delta_1+\delta_2)/2\ge\psi$.
Now $G_2\equiv D^2(D-1)$ implies $K\equiv D-1$.
In Laurent coefficients its central coefficient is
odd and at least one, its coefficients at $q^{\pm1}$
are odd and at least one, and those at $q^{\pm2}$
are even and at least two. Therefore $K(4)\ge7$,
so $t\ge14$. The preceding symmetric scalar estimate,
including the remaining lower bounds in (G031.1), gives
\[
 T\ge16t^2+458t+2696
    \ge12244>11128.
\]

\paragraph{The remaining pair $(4,4)$.}
Here (G031.5) becomes
\[
 g_1=\frac{12}{5}(3d_1+2d_2),\qquad
 g_2=\frac{12}{5}(2d_1+3d_2),\qquad
 g_1+g_2=12t.
\]
Also
\[
 s=\frac{24}{25}(6t^2+d_1d_2)\ge\frac{144}{25}t^2.
\]
Using the remaining mass bounds in (G031.1),
\[
 T\ge\frac{144}{25}t^2+266t+2336.                  \tag{G031.6}
\]
At $t=24$ this exceeds $11128$, and it is increasing
for $t\ge0$. Thus $t<24$.

Write $d_1=2a,d_2=2b$. The integrality of $g_1,g_2$
gives $b=a+5k$ for some integer $k$. The stronger
conditions $48\mid g_1$ and $g_2/24$ odd show
that $a$ is even and $k$ is odd:
\[
 g_1=24a+48k,\qquad g_2=24a+72k.
\]
With $a,b\ge2$ and $d_1+d_2<24$, the only possibilities
are $(d_1,d_2)=(4,14)$ and $(16,6)$.
For completeness, $k\ge1$ forces $t=4a+10k\ge18$;
if $k\ge3$ it is at least $38$, while $k=1$ forces
$a=2$. If $k\le-1$, the inequality $b\ge2$ gives
$a\ge2-5k$; evenness improves this to $a\ge8$ when
$k=-1$, leaving $a=8$, and $k\le-3$ forces $t\ge42$.

For $(4,14)$ one has $(g_1,g_2)=(96,120)$ and
$\delta_1=\psi$. In (G031.4), write $C_1=2H$;
then $H$ is integral and $H(4)=g_1/48=2$.
Since both $U_i=D$, the first identity gives
\[
 \delta_1=W C_1-D C_2=2W H-D C_2.
\]
As $D\mid\delta_1$ and $W(0)=2$, this forces
$H(0)=0$. Thus $D\mid H$ over $\mathbb Z[D]$, and
$H(4)$ is a multiple of four, contradicting $H(4)=2$.

For $(16,6)$ one has $(g_1,g_2,s)=(144,120,2880)$.
The remaining identity gives
\[
 v_1=\frac{576-a_1}{20}.
\]
Since $v_1$ is even, $a_1\equiv16\pmod{40}$;
with $a_1\ge8$ it follows that $a_1\ge16$.
Hence $14a_1+20v_1=576+13a_1\ge784$.
The other terms of (G031.1) then imply
\[
 T\ge2880+5280+1200+160+572+784+544
       =11420>11128.
\]
This excludes the last branch and proves the statement.
\hfill$\square$

\paragraph{Geometric translation and scope.}
In the NN-G010 realization by graded sign-representation
kernels, $w$ is the ordinary dimension of the gluing
coefficient object. The theorem says that any such
closed-table realization costing at most $3352$
regular-sign pairs must use a gluing object of dimension
at least eight. A dimension-six gluing object has the
unique character $D+2$; its associative kernel products
would force the incompatible character congruences and
mass bounds proved above. This is an actual dimension
obstruction for the effective realization, not a statement
about Frobenius weights or point counts. The character
dictionary is SGA 3, I, 4.7.2--4.7.3; the composition
dictionary is archive cote 2, scan 17, as realized in
NN-G010. The result shrinks the remaining optimization
gate but gives no global optimum or canonical GCT basis.

